% =====================================================================
%  The Honjo Masamune Language
%  A Cut-Primitive Domain-Specific Language for Cheminformatics by
%  Individuation
%
%  Language specification. Self-contained: the framework results it
%  binds to are stated as named operations, not re-derived.
% =====================================================================
\documentclass[11pt,a4paper]{article}

\usepackage[utf8]{inputenc}
\usepackage[T1]{fontenc}
\usepackage{lmodern}
\usepackage[a4paper,margin=2.4cm]{geometry}
\usepackage{amsmath,amssymb,amsthm,mathtools}
\usepackage{bm}
\usepackage{mathrsfs}
\usepackage{booktabs}
\usepackage{array}
\usepackage{enumitem}
\usepackage{microtype}
\usepackage{xcolor}
\usepackage{listings}
\usepackage[round,sort&compress]{natbib}
\usepackage[colorlinks=true,linkcolor=blue!45!black,
            citecolor=blue!45!black,urlcolor=blue!45!black]{hyperref}
\usepackage{cleveref}

% ---- Theorem-like environments ----
\newtheoremstyle{thmstyle}{\topsep}{\topsep}{\itshape}{}{\bfseries}{.}{.5em}{}
\theoremstyle{thmstyle}
\newtheorem{theorem}{Theorem}[section]
\newtheorem{lemma}[theorem]{Lemma}
\newtheorem{proposition}[theorem]{Proposition}
\newtheorem{corollary}[theorem]{Corollary}
\theoremstyle{definition}
\newtheorem{definition}[theorem]{Definition}
\newtheorem{rulename}[theorem]{Rule}
\newtheorem{principle}[theorem]{Principle}
\newtheorem{example}[theorem]{Example}
\theoremstyle{remark}
\newtheorem{remark}[theorem]{Remark}
\newtheorem{notation}[theorem]{Notation}

% ---- honjo listing style ----
\definecolor{hjkw}{RGB}{30,64,175}
\definecolor{hjcom}{RGB}{22,101,52}
\definecolor{hjnum}{RGB}{180,83,9}
\definecolor{hjstr}{RGB}{159,18,57}
\definecolor{hjbg}{RGB}{250,250,250}
\lstdefinelanguage{honjo}{
  morekeywords={floor,cut,close,track,until,yield,when,emit,observe,
    where,in,as,let,medium,residue,converge,diverge,with,by,reps,
    representation,switch,import,module,export,assert,from,to,if,else},
  morekeywords=[2]{atom,bond,compound,path,scalar,Z,nu,thick,delta,
    valence,shape,vacancy,floorval,count},
  sensitive=true,
  morecomment=[l]{--},
  morestring=[b]",
  keywordstyle=\color{hjkw}\bfseries,
  keywordstyle=[2]\color{hjnum},
  commentstyle=\color{hjcom}\itshape,
  stringstyle=\color{hjstr},
  basicstyle=\ttfamily\small,
  backgroundcolor=\color{hjbg},
  frame=single,framesep=4pt,xleftmargin=6pt,xrightmargin=4pt,
  breaklines=true,showstringspaces=false,columns=fullflexible,tabsize=2
}
\lstset{language=honjo}

\crefname{rulename}{Rule}{Rules}
\Crefname{rulename}{Rule}{Rules}

% ---- shorthand ----
\newcommand{\R}{\mathbb{R}}
\newcommand{\N}{\mathbb{N}}
\newcommand{\flo}{\beta}
\newcommand{\res}{\rho}
\newcommand{\cut}{\partial}
\newcommand{\medium}{\mathsf{m}}
\newcommand{\evalto}{\Downarrow}
\newcommand{\dd}{\,\mathrm{d}}
\newcommand{\bnf}{\;::=\;}
\newcommand{\alt}{\;\mid\;}
\newcommand{\hj}[1]{\texttt{#1}}
\newcommand{\abs}[1]{\lvert#1\rvert}

\title{\textbf{The Honjo Masamune Language:\\[3pt]
A Cut-Primitive Domain-Specific Language for\\
Cheminformatics by Individuation}}

\author{%
Kundai Farai Sachikonye\\
\small Technical University of Munich\\
\small AIMe Registry for Artificial Intelligence\\
\small \texttt{kundai.sachikonye@wzw.tum.de}}

\date{\today}

\begin{document}
\maketitle

% =====================================================================
\begin{abstract}
\noindent
We specify \emph{Honjo Masamune} (\hj{honjo}; source extension
\hj{.hj}), a domain-specific language for cheminformatics in which the
sole computational primitive is the \emph{cut}: the act of individuating
a part from the rest of a bounded whole at a strictly positive,
irreducible cost---the \emph{floor} $\flo>0$. The language is built on a
single semantic identity carried from its underlying theory: generating
an atom, forming a chemical bond, and tracking an item through a process
are not distinct operations but the same cut operation at different
arities. A length-one cut individuates an atom; a cut between two items is
a bond; a residue-chained sequence of cuts is a process track. honjo
therefore has one verb and one composition rule, from which atom
generation, compound construction, molecular geometry, and causal
process tracking all follow. We give: (i)~the semantic model---a finite
weighted \emph{contact graph} on which every honjo value is a cut of
accountable residue, never a bare number; (ii)~a complete lexical and
context-free grammar in EBNF; (iii)~a static \emph{accountability} type
system in which every value carries the floor at which it was cut, so
that a value claiming zero residue is rejected at compile time;
(iv)~a small-step operational semantics in which evaluation \emph{is}
measurement (executing a cut performs it; it does not simulate it), and
in which the committed-cut count is monotone, giving every program an
intrinsic, non-decreasing clock; (v)~a standard library of verbs, each
bound by name to a validated operation of the underlying framework
(\hj{cut} $\to$ spectroscopic atom individuation; \hj{$\sim$} and
\hj{close} $\to$ boundary-thickness bonding and shell closure;
\hj{track\ldots until\ldots yield} $\to$ accountable propagation with
$S$-entropy convergence admissibility); and (vi)~a two-target
compilation model. honjo compiles, from one front end and one typed core
intermediate representation, to two back ends: a TypeScript/WebGL\,2
runtime for the interactive online tool, in which a program's cuts are
realised as GPU render passes (rendering is measurement), and a Rust
runtime for the reference implementation, in which cuts are realised as
exact maximum-flow minimum-cut computations on the contact graph. We
prove that the two back ends are observationally equivalent up to the
floor: any honjo program yields the same accountable result on both,
differing only below the resolution the floor already forbids. The
specification is implementable as written; a companion reference
interpreter realises the operational semantics directly.
\end{abstract}

\noindent\textbf{Keywords:} domain-specific language; cheminformatics;
individuation; cut primitive; contact graph; accountability type system;
operational semantics; minimum cut; rendering-as-measurement; two-target
compilation; TypeScript; Rust.

\tableofcontents

% =====================================================================
\section{Introduction}
\label{sec:intro}
% =====================================================================

\subsection{Papers as programs}

The framework underlying this language consists of a sequence of results
that are, in form, \emph{operations}: a procedure that individuates an
atom from spectroscopic structure; a procedure that decides whether two
atoms bond and in what geometry; a procedure that tracks an item's
uncertainty through a process and returns an accountable account of what
happened to it. Each such result is already, in effect, a program with a
proof of correctness attached. This specification makes that literal: it
defines a language whose statements \emph{are} those operations, so that
a paper becomes a script and a derivation becomes a run.

The design discipline throughout is that we implement only what is
already established. Every verb of the language binds to an operation of
the framework that has an explicit construction and has been validated
numerically; the language adds composition, typing, and execution, not
new physics or chemistry.

\subsection{One primitive: the cut}

The language has a single computational primitive, the \emph{cut}, and a
single reason for it. To individuate any part of a bounded whole---to
make it a distinct thing---costs a strictly positive, irreducible
minimum, the floor $\flo>0$. This one fact is the engine of the
underlying theory, and it is the engine of the language. Every honjo
operation is a cut:

\begin{itemize}[leftmargin=1.4em,itemsep=2pt]
\item A \emph{generation} is a cut of arity one: individuating an atom of
atomic number $Z$ from the medium (\Cref{sec:stdlib-atom}).
\item A \emph{bond} is a cut of arity two: the single shared boundary
between two items, which exists iff joining lowers total boundary
thickness (\Cref{sec:stdlib-bond}).
\item A \emph{compound} is a closure of cuts: cuts continued until every
constituent's vacancy is driven to zero (\Cref{sec:stdlib-compound}).
\item A \emph{track} is a residue-chained sequence of cuts: the residue
of each cut is the cause of the next, and the chain is the propagation of
an item's uncertainty through a process (\Cref{sec:stdlib-track}).
\end{itemize}

Generation is not a different verb from tracking; it is a track of length
one. This is the central design decision and it is forced by the theory,
not chosen for convenience: in the underlying framework a reaction step
and a measurement step are the same cut, and generation, bonding, and
tracking are that cut at arities one, two, and many. honjo therefore has
one verb (\hj{cut}) and one composition rule (residue chaining); all
other surface forms are sugar over these.

\subsection{Two targets, one language}

honjo exists in two parallel implementations from a single specification:

\begin{enumerate}[leftmargin=1.6em,itemsep=2pt]
\item A \textbf{TypeScript} compiler and WebGL\,2 runtime for the
interactive online tool. Here a cut is realised as a GPU render pass: the
framebuffer that results from rendering \emph{is} the measurement, in
keeping with the rendering--measurement identity of the underlying
theory. This is the fast, exploratory, in-browser version.
\item A \textbf{Rust} compiler and native runtime for the reference and
production version. Here a cut is realised as an exact maximum-flow
minimum-cut computation on the contact graph. This is the
authoritative version against which the online tool is checked.
\end{enumerate}

Both share one front end (lexer, parser, type checker) and one typed
core intermediate representation (\Cref{sec:ir}); they differ only in the
back end that realises a cut. We prove (\Cref{thm:target-equiv}) that the
two back ends agree on every program up to the floor---the one resolution
at which the theory already forbids finer distinction---so the online
tool and the reference implementation never disagree about anything the
language can express.

\subsection{What this document specifies}

\Cref{sec:model} fixes the semantic model: the contact graph, the cut,
the floor, the residue, and accountability. \Cref{sec:lexical} gives the
lexical grammar; \Cref{sec:grammar} the context-free grammar in EBNF.
\Cref{sec:types} defines the accountability type system and proves its
soundness. \Cref{sec:opsem} gives the small-step operational semantics
and proves cut monotonicity and measurement-not-simulation.
\Cref{sec:stdlib} specifies the standard library, binding each verb to its
framework operation. \Cref{sec:ir} defines the core IR and the two-target
compilation model and proves target equivalence. \Cref{sec:examples}
gives worked programs. \Cref{sec:impl} states implementation requirements
for both targets. Throughout, framework operations are named and their
input/output contracts stated; their derivations are external to this
document.

% =====================================================================
\section{The semantic model}
\label{sec:model}
% =====================================================================

honjo is defined over a single mathematical object, the contact graph,
and a single operation on it, the cut. We fix these before any syntax.

\subsection{The contact graph}

\begin{definition}[Contact graph]
\label{def:contact-graph}
A \emph{contact graph} is a finite weighted graph $\Gamma=(V,E,w)$ with a
distinguished \emph{medium} vertex $\medium\in V$ adjacent to every other
vertex, and $w\colon E\to\R_{>0}$ a strictly positive weight. Vertices
$v\neq\medium$ are \emph{items}; an edge between two items is a
\emph{contact}; $w(e)$ is the cost of the separation it maintains. The
graph is the runtime state of a honjo program: every value the program
holds is a structure on $\Gamma$.
\end{definition}

\begin{definition}[Floor]
\label{def:floor}
The \emph{floor} of a contact graph is a constant $\flo>0$ with
$w(e)\ge\flo$ for every edge. No separation is free. The floor is a
program-level parameter (set by \hj{floor}, \Cref{sec:grammar}); its
positivity is invariant and enforced by the type system
(\Cref{sec:types}).
\end{definition}

\begin{definition}[Cut]
\label{def:cut}
A \emph{cut} of an item set $S$ ($\varnothing\neq S\subsetneq V$,
$\medium\notin S$) is the edge boundary
$\cut(S)=\{e\in E: \abs{e\cap S}=1\}$, of weight
$w(\cut(S))=\sum_{e\in\cut(S)}w(e)\ge\flo$. The cut individuates $S$ from
the rest. A cut \emph{event} commits one such boundary: it fixes the
edges of $\cut(S)$ as realised separations and advances the program's
cut count by one.
\end{definition}

\begin{definition}[Residue]
\label{def:residue}
The \emph{residue} of a cut event is the weight $\res=w(\cut(S))\ge\flo$
it deposits as committed boundary. The residue is the value the cut
produces and the cause of any subsequent cut: a chain of cuts is one in
which each residue is consumed as the boundary condition of the next
(\Cref{sec:stdlib-track}).
\end{definition}

\begin{principle}[Accountability]
\label{prin:accountability}
Every honjo value is a cut, and every cut carries its residue. There are
no bare numbers in honjo: a quantity is always paired with the floor at
which it was individuated. A value claiming residue below the floor, or
zero residue, is not a value---it is the forbidden sharp cut, and the
type system rejects it (\Cref{thm:no-zero-value}).
\end{principle}

\begin{remark}[Why a graph and not a number line]
\label{rem:why-graph}
honjo computes on a contact graph rather than on real numbers because the
quantities it manipulates---atomic identity, bond existence, process
accountability---are minimum cuts, not scalars. A bond's strength is the
weight of a shared boundary; an atom's identity is the invariant minimum
cut against the medium; a track's result is an accountable chain of cuts.
Representing these as graph cuts (rather than as fitted scalars) is what
lets the same primitive serve generation, bonding, and tracking, and is
what the two back ends each realise: exactly (Rust, max-flow) or by
rendering (TypeScript, GPU).
\end{remark}

\subsection{Measurement, not simulation}

\begin{principle}[Evaluation is measurement]
\label{prin:measure}
Executing a honjo cut performs the individuation; it does not simulate a
description of it. The result of \hj{cut\ 8} is the measured partition
structure of oxygen, carried as an accountable value, not a numerical
model of oxygen. Operationally (\Cref{sec:opsem}) this means a cut
evaluation mutates the contact graph---it commits boundary and advances
the count---rather than returning a pure function of inputs. The
distinction is observable: re-evaluating a cut is a new cut at a higher
count, never a cached recomputation (\Cref{thm:monotone}).
\end{principle}

% =====================================================================
\section{Lexical structure}
\label{sec:lexical}
% =====================================================================

A honjo source file is UTF-8 text with extension \hj{.hj}. Lexing
produces a token stream from the following classes.

\begin{definition}[Token classes]
\label{def:tokens}
\leavevmode
\begin{itemize}[leftmargin=1.4em,itemsep=1pt]
\item \textbf{Comments.} \hj{--} to end of line; discarded.
\item \textbf{Keywords.} \hj{floor}, \hj{cut}, \hj{close}, \hj{track},
\hj{until}, \hj{yield}, \hj{when}, \hj{do}, \hj{emit}, \hj{observe},
\hj{in}, \hj{as}, \hj{let}, \hj{medium}, \hj{converge}, \hj{diverge},
\hj{with}, \hj{by}, \hj{import}, \hj{module}, \hj{export}, \hj{assert}.
\item \textbf{Identifiers.} \hj{[A-Za-z\_][A-Za-z0-9\_]*}, not a keyword.
\item \textbf{Numbers.} decimal integers and floats with optional
exponent: \hj{[0-9]+(.[0-9]+)?([eE][+-]?[0-9]+)?}.
\item \textbf{Strings.} double-quoted, with the usual escapes.
\item \textbf{Operators and punctuation.} \hj{:=} (bind), \hj{\~{}} (cut
between / bond), \hj{( ) [ ] \{ \} , : . > < >= <= == ->} and the
range/bound forms used by \hj{cell}/\hj{bounds}.
\end{itemize}
Whitespace is insignificant except as a token separator; honjo is not
layout-sensitive.
\end{definition}

\begin{notation}[Floors and units in literals]
\label{not:literals}
A numeric literal may carry an explicit floor annotation with the suffix
form \hj{value\#floor} (read ``value at floor''); e.g.\ \hj{1.0e-7\#1e-9}.
A literal without an explicit floor inherits the ambient program floor set
by the most recent \hj{floor} declaration in scope. There is no way to
write a literal of zero floor (\Cref{thm:no-zero-value}).
\end{notation}

% =====================================================================
\section{Grammar}
\label{sec:grammar}
% =====================================================================

We give the context-free grammar in EBNF \citep{ebnf2017}. Nonterminals
are \textit{italicised}; terminals are \hj{typewriter}; \hj{\{x\}} is zero
or more, \hj{[x]} optional, \hj{|} alternation.

\begin{definition}[honjo grammar]
\label{def:grammar}
\begin{align*}
\textit{program} &\bnf \{\textit{decl}\} \\[2pt]
\textit{decl} &\bnf \textit{floorDecl} \alt \textit{import}
   \alt \textit{moduleDecl} \alt \textit{stmt} \\[2pt]
\textit{floorDecl} &\bnf \hj{floor}\ \textit{number} \\
\textit{import} &\bnf \hj{import}\ \textit{qname} \\
\textit{moduleDecl} &\bnf \hj{module}\ \textit{ident}\ \hj{\{}\ \{\textit{decl}\}\ \hj{\}} \\[2pt]
\textit{stmt} &\bnf \textit{bind} \alt \textit{track} \alt \textit{observe}
   \alt \textit{assert} \alt \textit{expr} \\[2pt]
\textit{bind} &\bnf [\hj{let}]\ \textit{ident}\ \hj{:=}\ \textit{expr} \\[2pt]
\textit{expr} &\bnf \textit{cut} \alt \textit{bondExpr}
   \alt \textit{closeExpr} \alt \textit{call} \alt \textit{ident}
   \alt \textit{number} \alt \textit{string}
   \alt \hj{(}\ \textit{expr}\ \hj{)} \\[2pt]
\textit{cut} &\bnf \hj{cut}\ \textit{cutArg} \\
\textit{cutArg} &\bnf \textit{number} \alt \textit{ident}
   \alt \textit{ident}\ \hj{\~{}}\ \textit{ident} \\[2pt]
\textit{bondExpr} &\bnf \textit{expr}\ \hj{\~{}}\ \textit{expr}\
   [\,\hj{when}\ \textit{cond}\,] \\[2pt]
\textit{closeExpr} &\bnf \hj{close}\ \textit{ident}\ \hj{(}\
   \textit{argList}\ \hj{)}\ [\,\hj{by}\ \textit{ident}\,] \\[2pt]
\textit{track} &\bnf \hj{track}\ \textit{ident}\ \hj{in}\ \textit{expr} \\
   &\qquad [\,\hj{with}\ \hj{reps}\ \textit{repList}\,]\ \\
   &\qquad \hj{until}\ \textit{admit}\ \ \hj{yield}\ \textit{ident} \\[2pt]
\textit{admit} &\bnf \hj{converge} \alt \hj{diverge} \alt \textit{cond} \\[2pt]
\textit{observe} &\bnf \hj{observe}\ \textit{expr}\
   [\,\hj{as}\ \textit{ident}\,] \\[2pt]
\textit{assert} &\bnf \hj{assert}\ \textit{cond}\
   [\,\hj{emit}\ \textit{string}\,] \\[2pt]
\textit{call} &\bnf \textit{qname}\ \hj{(}\ [\textit{argList}]\ \hj{)} \\
\textit{argList} &\bnf \textit{arg}\ \{\hj{,}\ \textit{arg}\} \\
\textit{arg} &\bnf [\textit{ident}\ \hj{:}]\ \textit{expr} \\[2pt]
\textit{cond} &\bnf \textit{expr}\ \textit{relop}\ \textit{expr} \\
\textit{relop} &\bnf \hj{>} \alt \hj{<} \alt \hj{>=} \alt \hj{<=}
   \alt \hj{==} \\[2pt]
\textit{repList} &\bnf \textit{ident}\ \{\hj{,}\ \textit{ident}\} \\
\textit{qname} &\bnf \textit{ident}\ \{\hj{.}\ \textit{ident}\} \\
\textit{number} &\bnf \textit{numlit}\ [\hj{\#}\ \textit{numlit}]
\end{align*}
\end{definition}

\begin{remark}[Reading the core forms]
\label{rem:core-forms}
The whole language is small. \hj{cut\ N} generates (arity-one cut);
\hj{a\ \~{}\ b\ when\ c} bonds (arity-two cut, guarded);
\hj{close\ X(\ldots)} drives an item to shell closure (cut-to-closure);
\hj{track\ x\ in\ P\ until\ converge\ yield\ r} propagates and binds the
accountable result. \hj{observe} forces measurement of a value (commits
its cut now); \hj{assert} is a pre-flight check that may halt with a
message. Everything else---\hj{let}, \hj{import}, \hj{module}, named
arguments---is conventional.
\end{remark}

% =====================================================================
\section{The accountability type system}
\label{sec:types}
% =====================================================================

honjo's types record \emph{what kind of cut} a value is and \emph{at what
floor}. The system's purpose is to make Principle~\ref{prin:accountability}
statically enforceable: a value of zero residue cannot be typed.

\subsection{Types}

\begin{definition}[Types]
\label{def:types}
The honjo types are
\[
\tau \bnf \mathsf{Atom} \alt \mathsf{Bond} \alt \mathsf{Compound}
   \alt \mathsf{Path} \alt \mathsf{Scalar} \alt \mathsf{Cut}.
\]
$\mathsf{Cut}$ is the supertype of all individuated values; $\mathsf{Atom}$,
$\mathsf{Bond}$, $\mathsf{Compound}$, $\mathsf{Path}$ are its arity- and
role-refinements ($\mathsf{Atom}=$ arity-one cut, $\mathsf{Bond}=$
arity-two cut, $\mathsf{Compound}=$ closure of cuts, $\mathsf{Path}=$
residue chain). $\mathsf{Scalar}$ is a measured number; it is never bare.
Every typed value $v$ additionally carries a \emph{floor annotation}
$\flo(v)>0$ and a \emph{residue} $\res(v)\ge\flo(v)$.
\end{definition}

\begin{definition}[Typing judgement]
\label{def:judgement}
A judgement $\Gamma\vdash e : \tau \,@\, \flo$ reads ``in floor context
$\Gamma$ (which fixes the ambient floor), expression $e$ has type $\tau$
and is individuated at floor $\flo>0$.'' The context carries the ambient
floor set by the most recent \hj{floor} declaration and the types of
bound identifiers.
\end{definition}

\subsection{The positivity rule}

\begin{rulename}[Floor positivity]
\label{rule:positive}
Every value-introducing rule has the side condition $\flo>0$ and
$\res\ge\flo$. A literal or expression whose annotation forces
$\flo\le0$ or $\res<\flo$ is ill-typed.
\[
\frac{\Gamma\vdash \textit{numlit}\;n \quad \flo>0}
     {\Gamma\vdash n\#\flo : \mathsf{Scalar}\,@\,\flo}
\qquad
\frac{}{\Gamma\vdash \hj{cut}\ Z : \mathsf{Atom}\,@\,\flo_\Gamma}
\]
where $\flo_\Gamma$ is the ambient floor.
\end{rulename}

\begin{theorem}[No zero-residue value]
\label{thm:no-zero-value}
There is no well-typed honjo expression of residue $0$. Equivalently, the
sharp cut is not expressible: every typeable value has $\res\ge\flo>0$.
\end{theorem}

\begin{proof}
By induction on typing derivations. The only value introductions are
literals (\Cref{rule:positive}, side condition $\flo>0$, and a literal's
residue is its annotated floor), cuts (residue $=w(\cut(S))\ge\flo>0$ by
\Cref{def:cut}), and the derived forms (bond, close, track) whose residues
are sums or chains of cut residues, hence $\ge\flo$. No rule introduces a
value of residue $<\flo$, and $\flo>0$ is a side condition on every
floor-introduction. Hence no derivation concludes $\res=0$.
\end{proof}

\subsection{Cut typing}

\begin{rulename}[Generation; bond; close]
\label{rule:verbs}
\[
\frac{\Gamma\vdash Z:\mathsf{Scalar}\quad Z\in\N^{+}}
     {\Gamma\vdash \hj{cut}\ Z : \mathsf{Atom}\,@\,\flo_\Gamma}
\qquad
\frac{\Gamma\vdash a:\mathsf{Atom}\quad \Gamma\vdash b:\mathsf{Atom}}
     {\Gamma\vdash a\ \hj{\~{}}\ b : \mathsf{Bond}\,@\,\flo_\Gamma}
\]
\[
\frac{\Gamma\vdash X:\mathsf{Atom}\quad
      \Gamma\vdash a_i:\mathsf{Atom}}
     {\Gamma\vdash \hj{close}\ X(a_1,\dots,a_k) : \mathsf{Compound}\,@\,\flo_\Gamma}
\]
The bond rule additionally requires its guard, when present, to be a
$\mathsf{Bool}$-valued $\textit{cond}$; an unguarded bond is admitted only
if static analysis cannot exclude $\Delta\mathrm{thick}>0$
(\Cref{sec:stdlib-bond}).
\end{rulename}

\begin{rulename}[Track]
\label{rule:track}
\[
\frac{\Gamma\vdash x:\mathsf{Atom}\ \text{or}\ \mathsf{Compound}\quad
      \Gamma\vdash P:\mathsf{Compound}\ \text{or}\ \mathsf{Path}}
     {\Gamma\vdash \hj{track}\ x\ \hj{in}\ P\ \hj{until}\ \textit{admit}\ \hj{yield}\ r
      : \mathsf{Path}\,@\,\flo_\Gamma}
\]
The bound result $r$ has type $\mathsf{Path}$; its residue is the total
committed boundary of the propagation (\Cref{sec:stdlib-track}).
\end{rulename}

\begin{theorem}[Type soundness]
\label{thm:soundness}
honjo typing is sound for the operational semantics of
\Cref{sec:opsem}: if $\Gamma\vdash e:\tau\,@\,\flo$ and $e$ evaluates to a
value $v$, then $v$ has type $\tau$, floor $\flo>0$, and residue
$\res(v)\ge\flo$ (progress and preservation).
\end{theorem}

\begin{proof}[Proof sketch]
Standard syntactic soundness \citep{pierce2002}. \emph{Progress}: a
well-typed closed expression is a value or steps (every form has a
reduction rule in \Cref{sec:opsem}). \emph{Preservation}: each reduction
rule preserves the type and the floor annotation, and---by
\Cref{thm:no-zero-value} applied at each step---preserves $\res\ge\flo>0$.
The only non-standard clause is that reductions are state-mutating
(measurement, \Cref{prin:measure}); preservation holds because the cut
count and committed boundary are monotone (\Cref{thm:monotone}), so no
step lowers a residue below the floor.
\end{proof}

% =====================================================================
\section{Operational semantics}
\label{sec:opsem}
% =====================================================================

We give a small-step semantics over configurations
$\langle \Gamma_G, M, e\rangle$, where $\Gamma_G$ is the current contact
graph (\Cref{def:contact-graph}), $M\in\N$ is the committed-cut count
(the program clock), and $e$ is the expression under evaluation. We write
$\langle\Gamma_G,M,e\rangle \to \langle\Gamma_G',M',v\rangle$ for one
step. Evaluation is measurement: steps mutate $\Gamma_G$ and increment
$M$.

\subsection{Cut reduction}

\begin{rulename}[Cut commits and counts]
\label{rule:cut-step}
\[
\frac{S=\textsc{Individuate}(\Gamma_G,Z)\qquad
      \res=w(\cut(S))\ge\flo}
     {\langle\Gamma_G,M,\hj{cut}\ Z\rangle \to
      \langle\Gamma_G\!\oplus\!\cut(S),\,M+1,\,
      v_{\mathsf{Atom}}(S,\res)\rangle}
\]
where $\textsc{Individuate}$ is the framework's atom-individuation
operation (\Cref{sec:stdlib-atom}), $\Gamma_G\!\oplus\!\cut(S)$ is
$\Gamma_G$ with the cut's boundary committed, and $v_{\mathsf{Atom}}(S,\res)$
is the resulting accountable atom value.
\end{rulename}

The bond, close, and track rules are analogous: each commits boundary and
increments $M$ by the number of cut events it performs (one for a bond,
the closure count for \hj{close}, the chain length for \hj{track}).

\begin{theorem}[Cut monotonicity / intrinsic clock]
\label{thm:monotone}
For any program, the committed-cut count $M$ is non-decreasing along
evaluation and strictly increases at every cut event. No reduction
decreases $M$. Consequently re-evaluating the same surface expression is a
new cut at a strictly higher count, never a cached recomputation, and two
configurations with identical graphs but different $M$ are distinct
states.
\end{theorem}

\begin{proof}
Every value-producing rule (\Cref{rule:cut-step} and its analogues)
increments $M$; no rule decrements it (there is no reduction that
un-commits boundary---``undo'' is itself a further cut). Hence $M$ is
monotone and strictly increases per cut event. Distinctness of equal-graph
states at different $M$ is immediate since $M$ is part of the
configuration.
\end{proof}

\begin{corollary}[Measurement, not simulation]
\label{cor:measure}
Because evaluation mutates $\Gamma_G$ and advances $M$, a honjo cut cannot
be a pure function memoised across calls: each occurrence performs the
individuation afresh and is recorded in the clock. This realises
Principle~\ref{prin:measure}: running the program \emph{is} performing the
measurements it names.
\end{corollary}

\subsection{Track reduction and convergence}

\begin{rulename}[Track is a residue chain]
\label{rule:track-step}
\[
\frac{\begin{array}{c}
       (S_1,\dots,S_k)=\textsc{Propagate}(\Gamma_G,x,P)\quad
       \res_i=w(\cut(S_i))\\
       \textsc{Admit}(\res_1,\dots,\res_k)=\textit{admit}
      \end{array}}
     {\langle\Gamma_G,M,\hj{track}\ x\ \hj{in}\ P\ \hj{until}\ \textit{admit}\ \hj{yield}\ r\rangle
      \to \langle\Gamma_G',\,M+k,\,
      v_{\mathsf{Path}}(S_{1:k},\textstyle\sum_i\res_i)\rangle}
\]
$\textsc{Propagate}$ is the framework's accountable-propagation operation;
each $\res_i$ is the residue of cut $i$ and the boundary condition of cut
$i{+}1$ (residue chaining). $\textsc{Admit}$ tests the convergence
criterion ($S$-entropy convergence for \hj{converge}); a non-admissible
chain reduces to the empty path (no yield).
\end{rulename}

\begin{remark}[Local steps are free; only convergence is judged]
\label{rem:local-free}
In keeping with the underlying propagation theory, the intermediate cuts
of a track are unconstrained: a step may be locally extreme (even
``impossible'') provided the chain composes to a value that meets the
admissibility criterion. \textsc{Admit} therefore inspects only the
terminal convergence of the chain, not the plausibility of intermediate
$\res_i$. This is what makes \hj{track\ldots until\ converge} a search over
admissible propagations rather than a single forced path.
\end{remark}

% =====================================================================
\section{The standard library}
\label{sec:stdlib}
% =====================================================================

Each honjo verb binds to a named operation of the underlying framework.
We state the contract (inputs, output, the operation it invokes); the
operation's construction and validation are external to this document.

\subsection{\texorpdfstring{\hj{cut} --- atom individuation}{cut --- atom individuation}}
\label{sec:stdlib-atom}

\begin{definition}[\textsc{Individuate}]
\label{def:individuate}
\hj{cut\ Z} invokes \textsc{Individuate}$(Z)$: from atomic number
$Z\in\N^{+}$, fill the partition coordinates $(n,\ell,m,s)$ under the
shell capacity $C(n)=2n^2$, exclusion, and energy ordering, returning an
$\mathsf{Atom}$ value carrying the configuration, the valence-shell
occupancy $q_v$ and capacity $C_v$, the vacancy $\nu=C_v-q_v$, the ground
term symbol, and the residue $\res=w(\cut(\{Z\}))\ge\flo$. This is the
spectroscopic atom-derivation operation; e.g.\ \hj{cut\ 6} yields carbon
$[\mathrm{He}]2s^22p^2$, $^3P_0$, $\nu=4$.
\end{definition}

\subsection{\texorpdfstring{\hj{$\sim$} --- bonding}{~ --- bonding}}
\label{sec:stdlib-bond}

\begin{definition}[\textsc{Bond}]
\label{def:bond}
\hj{a\ \~{}\ b} invokes \textsc{Bond}$(a,b)$: form the shared-boundary cut
between atoms $a,b$ and return a $\mathsf{Bond}$ value iff the shared
content $\Delta\mathrm{thick}(a,b)>0$ (joint individuation is cheaper than
separate), bounded by $\min(\nu_a,\nu_b)$; otherwise no bond. The optional
guard \hj{when\ }$c$ requires $c$ to hold (e.g.\ \hj{when\ delta\ >\ 0}).
This is the boundary-thickness bonding criterion; a closed-shell atom
($\nu=0$) bonds with nothing.
\end{definition}

\subsection{\texorpdfstring{\hj{close} --- compound by closure}{close --- compound by closure}}
\label{sec:stdlib-compound}

\begin{definition}[\textsc{Close}]
\label{def:close}
\hj{close\ X(a$_1$,\ldots,a$_k$)} invokes \textsc{Close}: form the set of
bonds that drives every constituent's vacancy to zero, returning a
$\mathsf{Compound}$ value carrying the stoichiometry (fixed by vacancy
matching), the per-centre geometry (the maximal angular separation of the
committed interfaces, e.g.\ $109.47^\circ$ tetrahedral, compressed by
unshared regions), and total residue. E.g.\ \hj{close\ O(H,H)} yields a
bent $2{:}1$ water at $\approx104.5^\circ$.
\end{definition}

\subsection{\texorpdfstring{\hj{track} --- accountable propagation}{track --- accountable propagation}}
\label{sec:stdlib-track}

\begin{definition}[\textsc{Propagate}, \textsc{Admit}]
\label{def:propagate}
\hj{track\ x\ in\ P\ until\ converge\ yield\ r} invokes
\textsc{Propagate}$(x,P)$: track item $x$'s individuation through process
$P$ as a residue-chained sequence of cuts, and \textsc{Admit} the chain
iff it $S$-entropy-converges to the output (the only admissibility test;
local steps free, \Cref{rem:local-free}). The yielded $\mathsf{Path}$ is
the \emph{amalgamation}: the accountable set of cuts---reaction steps and
measurement steps indifferently---by which $x$ reached the output.
Optional \hj{with\ reps\ }$R$ permits representation switching between
steps (mean-recovery-equivalent decompositions) at no additional cut cost.
\end{definition}

\begin{remark}[The verbs are one verb]
\label{rem:one-verb}
\textsc{Individuate}, \textsc{Bond}, \textsc{Close}, and
\textsc{Propagate} are the cut at arities one, two, closure, and chain.
The standard library is therefore not four independent operations but one
operation (the cut) exposed at four arities, exactly as
\Cref{sec:intro} claims and as the underlying theory requires.
\end{remark}

% =====================================================================
\section{Core IR and the two-target model}
\label{sec:ir}
% =====================================================================

honjo compiles, from one front end, to a typed core intermediate
representation, and thence to two back ends. This section fixes the IR and
proves the two back ends agree.

\subsection{The core IR}

\begin{definition}[Cut-IR]
\label{def:ir}
\emph{Cut-IR} is a typed sequence of cut instructions over a contact-graph
state. Each instruction is one of: $\textsc{Ind}(Z)$ (individuate),
$\textsc{Bnd}(a,b,g)$ (guarded bond), $\textsc{Cls}(X,\bar a)$ (close),
$\textsc{Prp}(x,P,\textit{adm})$ (propagate), $\textsc{Obs}(v)$ (force
measurement). Each instruction carries its result type and floor
(\Cref{sec:types}) and, when executed, returns an accountable value and
increments the cut count. Cut-IR is back-end-agnostic: it names cuts, not
how they are realised.
\end{definition}

The front end (lexer, parser, type checker of \Cref{sec:lexical}--%
\ref{sec:types}) is shared by both targets and lowers a program to Cut-IR.
Only the realisation of a cut instruction differs between targets.

\subsection{Target R (Rust): cuts as minimum cuts}

\begin{definition}[Rust realisation]
\label{def:rust-target}
In the Rust back end, each Cut-IR instruction is realised by an exact
computation on the contact graph: $\textsc{Ind}$, $\textsc{Bnd}$,
$\textsc{Cls}$ compute the relevant minimum cut by maximum
flow~\citep{fordfulkerson1956,cormen2009}; $\textsc{Prp}$ enumerates and
evaluates admissible propagations exactly; residues are exact edge-weight
sums. This is the reference semantics: it computes
$w(\cut(\cdot))$ to the floating-point precision of the host.
\end{definition}

\subsection{Target T (TypeScript/WebGL): cuts as render passes}

\begin{definition}[TypeScript realisation]
\label{def:ts-target}
In the TypeScript/WebGL\,2 back end, each Cut-IR instruction is realised
as a GPU render pass over the contact-graph state encoded in textures: a
fragment shader evaluates the cut at each texel and writes the
framebuffer, whose contents \emph{are} the accountable result
(rendering--measurement identity). Minimum cuts are computed by the
GPU-resident flow/observation kernels; residues are read back from the
framebuffer. This is the interactive online realisation.
\end{definition}

\subsection{Target equivalence}

\begin{theorem}[Two-target equivalence up to the floor]
\label{thm:target-equiv}
Let $p$ be a well-typed honjo program with ambient floor $\flo$. Let
$v_R$ and $v_T$ be the values it yields under the Rust and
TypeScript back ends respectively. Then $v_R$ and $v_T$ have the same
type and the same cut structure, and their residues agree up to the
floor:
\[
\abs{\res(v_R)-\res(v_T)}\;\le\;\flo .
\]
The two back ends never disagree about any distinction honjo can express,
because honjo cannot express a distinction finer than $\flo$
(\Cref{thm:no-zero-value}).
\end{theorem}

\begin{proof}
Both back ends realise the same Cut-IR (\Cref{def:ir}), so they perform
the same sequence of cut instructions with the same types and the same
cut count (the front end is shared). They differ only in how each cut's
residue is computed: Rust by exact max-flow, TypeScript by GPU
render-readback. The GPU realisation computes the same minimum cut to its
texel/precision resolution; its readback error is bounded by the
coarsest resolvable boundary, which is the floor (no honjo value is
individuated below $\flo$, \Cref{thm:no-zero-value}). Hence the residues
agree to within $\flo$, and since honjo admits no sub-floor distinction,
the two results are observationally identical at the language level.
\end{proof}

\begin{remark}[Why two targets, not one]
\label{rem:two-targets}
The split is deliberate and the equivalence theorem is what makes it
safe. The TypeScript/WebGL target is fast, in-browser, and renders cuts as
measurements directly---ideal for the interactive sandbox. The Rust target
is exact and authoritative---ideal as the reference against which the
online tool is validated. \Cref{thm:target-equiv} guarantees the online
tool is never \emph{wrong} relative to the reference; it is at most less
precise, and only below the floor, where the theory forbids meaning
anyway.
\end{remark}

% =====================================================================
\section{Worked programs}
\label{sec:examples}
% =====================================================================

\subsection{Generate an atom}

\begin{lstlisting}[caption={Individuating carbon (a length-one cut).}]
-- carbon.hj
floor 1.0                 -- set the ambient floor; nothing is free

C := cut 6                -- individuate Z=6 against the medium
observe C                 -- force the measurement now
-- C : Atom @ 1.0
--   config  [He] 2s2 2p2
--   term    3P_0
--   vacancy 4
\end{lstlisting}

\subsection{Build a compound}

\begin{lstlisting}[caption={Water: cuts to closure, geometry by separation.}]
-- water.hj
floor 1.0

O := cut 8
H := cut 1

-- a bond is a cut between two items, admitted only if it lowers thickness
OH := O ~ H when delta > 0

-- close drives every vacancy to zero: stoichiometry + geometry follow
W  := close O(H, H)       -- 2:1, bent, ~104.5 deg
observe W
assert W.valence == closed  emit "water did not close"
\end{lstlisting}

\subsection{Run the causal table}

\begin{lstlisting}[caption={Tracking an item through a process.}]
-- track.hj
floor 1.0
import honjo.causal

O := cut 8
W := close O(H, H)

-- propagate O's uncertainty through the process; admissible iff it
-- S-entropy-converges to the observed output
path := track O in W
          with reps mass, charge, time   -- representation switching allowed
          until converge
          yield amalgamation

observe path              -- the amalgamation IS the result
\end{lstlisting}

\begin{remark}[The same keyword, three jobs]
\label{rem:same-keyword}
Across the three programs, \hj{cut} generates, \hj{\~{}} (a cut between)
bonds, \hj{close} (cuts to closure) builds, and \hj{track} (a residue
chain of cuts) propagates---all the same primitive at different arities
(\Cref{rem:one-verb}). A reader who understands the cut understands the
whole language.
\end{remark}

% =====================================================================
\section{Implementation requirements}
\label{sec:impl}
% =====================================================================

Both targets must implement, from this specification:

\begin{enumerate}[leftmargin=1.6em,itemsep=2pt]
\item The shared front end: lexer (\Cref{sec:lexical}), parser to the
grammar (\Cref{def:grammar}), and accountability type checker
(\Cref{sec:types}) enforcing \Cref{thm:no-zero-value}.
\item Lowering to Cut-IR (\Cref{def:ir}).
\item The operational semantics of \Cref{sec:opsem}: a contact-graph
state, a monotone cut count, and the cut/bond/close/track reductions, with
evaluation mutating state (measurement, \Cref{cor:measure}).
\item The standard-library bindings of \Cref{sec:stdlib} to the framework
operations \textsc{Individuate}, \textsc{Bond}, \textsc{Close},
\textsc{Propagate}.
\item A back end: exact max-flow minimum cut (Rust,
\Cref{def:rust-target}) or GPU render-pass realisation (TypeScript/WebGL,
\Cref{def:ts-target}).
\end{enumerate}

A conformance suite must check \Cref{thm:target-equiv}: every example
program yields type- and structure-identical results on both targets with
residues agreeing to within the floor.

\begin{remark}[Reference interpreter]
\label{rem:reference}
A direct interpreter of \Cref{sec:opsem} over the exact (max-flow) back
end constitutes the simplest conforming implementation and serves as the
oracle for both production targets. It can reuse the already-validated
framework operations directly as the standard-library bindings.
\end{remark}

% =====================================================================
\section{Conclusion}
\label{sec:conclusion}
% =====================================================================

We have specified Honjo Masamune as a language whose single primitive is
the cut---the act of individuating a part of a bounded whole at the
irreducible cost $\flo>0$---and whose every operation, from generating an
atom to tracking an item through a process, is that primitive at a
different arity. The accountability type system makes the floor a static
invariant: no honjo value can claim the forbidden sharp cut. The
operational semantics makes evaluation measurement rather than
simulation, with a monotone cut count as the program's intrinsic clock.
The standard library binds each verb to a validated operation of the
underlying framework, so a honjo program is an executable paper. And the
two-target model---one shared, typed front end lowering to a
back-end-agnostic Cut-IR, realised exactly in Rust and by GPU rendering in
TypeScript---gives both an authoritative reference and an interactive
online tool that, by \Cref{thm:target-equiv}, never disagree about
anything the language can express. The specification is complete and
implementable as written; the reference interpreter of \Cref{rem:reference}
realises it directly from the validated framework operations.

\bibliographystyle{plainnat}
\bibliography{references}

\end{document}